\subsection{Positive cyclotomic eigenvalues}
\label{subsec:positive-eigenvalues}

The arithmetic cohomology also determines which scalar
denominators can vanish at the odd moment weights.
Put $P=\operatorname{Gal}(\Omega_p/F_1)$ and write
$u=\kappa(\gamma)$ for the cyclotomic value of a generator
of $\Gamma$.

\begin{lemma}\label{lem:positive-eigenvalues}
For every integer $m\geq2$,
\[
 \ker(\gamma-u^m:H^{\rm ab}\longrightarrow H^{\rm ab})=0.
\]
Consequently the characteristic polynomial of $\gamma$
on $D\oplus T$, or on any of its rational subquotients,
has nonzero value at $u^m$.
\end{lemma}
\begin{proof}
The freeness theorem for $H$ cited in Section~\ref{sec:abelian}
gives $\operatorname{cd}_pH\leq1$; also
$\operatorname{cd}_p\Gamma=1$.
For $A_n=\mathbb Z/p^n(m)$, Hochschild--Serre for
$1\to H\to P\to\Gamma\to1$ gives the edge isomorphism
\[
 H^2(P,A_n)\simeq
 \frac{\operatorname{Hom}_{\rm cts}(M,A_n)}
      {(\gamma-1)\operatorname{Hom}_{\rm cts}(M,A_n)}.
\]
The action is $(\gamma f)(x)=u^m f(\gamma^{-1}x)$.
Choose a coefficient generator and pair this discrete
module with $M/p^nM$. The adjoint of $\gamma-1$ is
$u^m\gamma^{-1}-1$. Thus its Pontryagin dual is exactly
\begin{equation}\label{eq:positive-eigenvalue-dual}
 H^2(P,A_n)^\vee\simeq
       \ker(\gamma-u^m:M/p^nM\longrightarrow M/p^nM).
\end{equation}
The choice rescales the pairing by a unit and changes
neither this kernel nor its order.

Let $\mathcal G=G_{F_1,S}$, with $S$ the primes above $p$
and the archimedean places, and put
$J=\ker(\mathcal G\to P)$. The coefficients $A_n$ factor
through $P$. A $P$-equivariant homomorphism $f:J\to A_n$
would give the exact extension
\[
 1\longrightarrow\operatorname{im}f\longrightarrow
       \mathcal G/\ker f\longrightarrow P\longrightarrow1.
\]
Its middle group is pro-$p$, so maximality of $P$ forces
$f=0$. Hence $H^1(J,A_n)^P=0$, and the five-term sequence
gives an injection
\[
                    H^2(P,A_n)\lhook\joinrel\longrightarrow
                         H^2(\mathcal G,A_n).
\]

The Chern-class comparison and even $K$-group finiteness
\cite[VI, Theorem~8.2, chapter p.~47]{Weibel} give the
finite group
\[
 H^2(\mathcal G,\mathbb Z_p(m))
     \simeq K_{2m-2}(O_{F_1}[1/p])\otimes\mathbb Z_p.
\]
Here $m\geq2$ is essential. Since $\mathcal G$ has
$p$-cohomological dimension two, continuous
$H^3(\mathcal G,\mathbb Z_p(m))=0$: the finite-coefficient
degree-three groups vanish, and the inverse system of
finite degree-two groups has zero $\varprojlim^1$.
The coefficient exact sequence therefore identifies
\[
 H^2(\mathcal G,A_n)
      =H^2(\mathcal G,\mathbb Z_p(m))/p^n.
\]
Its order, and hence the order in
\eqref{eq:positive-eigenvalue-dual}, is bounded independently of $n$.

A nonzero $u^m$-eigenvector in $M$ can be divided by its
maximal power of $p$ to give one outside $pM$.
Indeed $M$ is a product of copies of $\mathbb Z_p$ and
has no $p$-torsion, since it is the abelianization of
a free pro-$p$ group. Its image modulo $p^n$ would then
have order $p^n$ in the kernel in
\eqref{eq:positive-eigenvalue-dual}, contradicting the
uniform bound. This proves the first assertion.
A rational eigenvector in $D\oplus T$ can have its
denominator cleared in the integral lattice, so none
exists. The characteristic polynomials of rational
subquotients divide that of the whole finite-dimensional
module, proving the last assertion.
\end{proof}

Every nonempty product of positive odd target weights
has eigenvalue $u^{m_1+\cdots+m_d}$ with each $m_i\geq3$.
The lemma therefore shows that all companion denominators
coming from $D\oplus T$ are nonzero on this entire
moment grid. It also gives the exact cohomological reason
for this nonvanishing; the construction by a cofinite
grid in the preceding proof gives the same quotient map.
